% CV LaTeX — Ouahabi Benhenni
\documentclass[11pt,a4paper,sans]{moderncv}

% Style et couleurs
\moderncvstyle{banking}
\moderncvcolor{grey}

% Encodage et langue
\usepackage[T1]{fontenc}
\usepackage[utf8]{inputenc}
\usepackage[french]{babel}
\usepackage{microtype}

% Forcer les icônes symbols pour éviter fontawesome5
\moderncvicons{symbols}

% Mise en page
\usepackage[scale=0.90]{geometry}
\setlength{\hintscolumnwidth}{3.1cm}
\renewcommand*{\namefont}{\fontsize{26}{28}\bfseries\upshape}
\renewcommand*{\titlefont}{\fontsize{12}{14}\mdseries\upshape}

% Informations personnelles
\name{Ouahabi}{Benhenni}
\title{AI Software Engineer · Computer Vision \& Biomedical AI}
\address{Boumerdès, Algérie}{}{ }
\phone[mobile]{+213 6 76 95 79 17}
\email{Ouahabi.benhenni.fs@univ-boumerdes.dz}
\homepage{ouahabi-benhenni.web.app}
\social[linkedin]{ouahabi-benhenni-a28bb5155}
\social[github]{Greyma}

\begin{document}
\makecvtitle

% RÉCOMPENSES — mis en avant dès le début
\section{Récompenses et Distinctions}
\cvitem{2025}{\textbf{1\textsuperscript{ère} place} — Compétition Nationale IA Algérie-Tunisie (DevUp Boumerdès / Univ.\ de Tébessa) : projet de vision médicale par Deep Learning.}

% PROFIL
\section{Profil}
\cvitem{}{AI Software Engineer spécialisé en Computer Vision, Biomedical AI et systèmes distribués. Maîtrise des pipelines complets : acquisition de données, entraînement de modèles (CNN, Vision Transformers, U-Net, YOLOv8), déploiement et intégration dans des applications web temps réel. Recherche active en génomique computationnelle appliquée à l'identification de traitements via l'analyse ADN. Expérimenté en architecture microservices, encadrement technique et collaboration en équipe Agile.}

% COMPÉTENCES
\section{Compétences}
\cvitem{Langages}{Python, JavaScript, Java, C++, C\#, SQL, PHP}
\cvitem{IA / ML / DL}{PyTorch, TensorFlow, Keras, Scikit-Learn, OpenCV, Ultralytics, HuggingFace Transformers}
\cvitem{Architectures}{CNN, U-Net, YOLOv8, Vision Transformers, Mask R-CNN, RNN, Nucleotide Transformer}
\cvitem{LLM / GenAI}{LangChain, LangGraph, agents LLM, RAG, FAISS, ONNX}
\cvitem{Backend}{FastAPI, Flask, Django (DRF, Channels, Celery), Node.js (Express)}
\cvitem{Frontend}{React.js, Next.js, Tailwind CSS, ShadCN}
\cvitem{Bases de données}{PostgreSQL, MySQL, MongoDB, SQL Server}
\cvitem{DevOps \& Cloud}{Docker, CI/CD (GitHub Actions), Linux/VPS, Apache, Nginx, PM2}
\cvitem{Méthodologies}{OOP, Design Patterns, Git/GitHub, Agile/Scrum, RBAC, JWT}

% EXPÉRIENCE
\section{Expérience Professionnelle}

\cventry{Août 2025 -- Présent}{Consultant en Intelligence Artificielle}{Badis.ai}{Alger, Algérie}{}{%
\begin{itemize}
  \item \textbf{Audit \& Conseil IA} : Évaluation de faisabilité et planification d'un système d'IA adaptative éducative pour la plateforme ActivMath (1\,000+ utilisateurs).
  \item \textbf{Architecture} : Conception du pipeline de recommandation personnalisée et de suivi d'apprentissage (microservices Python, APIs REST, ONNX).
  \item \textbf{Intégration} : Déploiement de la solution dans l'écosystème existant (Flask/Next.js) sur VPS auto-géré.
\end{itemize}}

\cventry{Juin -- Septembre 2025}{Développeur Full-Stack}{Kleerinfinit}{Alger, Algérie}{}{%
\begin{itemize}
  \item Développement d'une plateforme SaaS B2B multi-tenant avec Django 5, DRF, Django Channels (WebSockets) et Celery.
  \item Conception d'une API RESTful avec authentification JWT, gestion des rôles RBAC et tableaux de bord analytiques.
  \item Collaboration en équipe de 5 développeurs (Agile, code reviews, sprints bimensuels).
\end{itemize}}

\cventry{2025}{Mentor — Bootcamp \og{}CodeCraft AI to FullStack\fg{}}{FormationTech}{Algérie}{}{%
\begin{itemize}
  \item Encadrement de 30+ participants sur Data Science, Deep Learning, Computer Vision et NLP.
  \item Supervision de projets de 48h présentés devant jury universitaire ; déploiement d'une plateforme de ressources IA.
\end{itemize}}

\cventry{2024--2025}{Développeur IA Freelance}{}{}{}{%
\begin{itemize}
  \item Conception de pipelines de segmentation et classification d'images médicales (U-Net, YOLOv8, LinkNet, PSPNet).
  \item Intégration des modèles dans des interfaces web interactives (Flask, React.js).
\end{itemize}}

\cventry{2025}{Chef du Département Développement — Club DevUp}{Université de Boumerdès}{}{}{%
Encadrement de 120 membres ; organisation des formations Web, Mobile et IA ; gestion de projets de développement.}

\cventry{2024}{Président du Club Informatique}{Université de Boumerdès}{}{}{%
Direction d'un club de 300 membres : hackathons, ateliers techniques, formations et partenariats universitaires.}

\cventry{2023--2024}{Projets IA Médicale}{Université de Boumerdès / CHU Mustapha Bacha}{}{}{%
\begin{itemize}
  \item Classification du cancer de la peau (CNN + Vision Transformers) — \textbf{93\,\% de précision}.
  \item Détection et comptage cellulaire automatisé (YOLOv8) ; détection du paludisme.
  \item Pipeline de diagnostic médical complet : preprocessing → inférence → interface web.
\end{itemize}}

% RECHERCHE
\section{Recherche \& Publications}
\cvitem{Mémoire M2}{\emph{Identification d'un traitement via ADN} — fine-tuning de Nucleotide Transformer v2 sur variants ClinVar pour la prédiction de pathogénicité SNV (AUC 0,987 en évaluation OOD par gène).}
\cvitem{Article soumis}{Classification des mélanomes par Deep Learning (CNN \& Vision Transformers) — soumis à comité de lecture.}

% FORMATION
\section{Formation}
\cventry{2026 (prévue)}{Master 2 Informatique — Systèmes d'Information}{Université de Boumerdès (UMBB)}{}{}{}
\cventry{2024}{Licence en Informatique — Systèmes d'Information}{Université de Boumerdès (UMBB)}{}{}{}
\cventry{2018--2021}{Technicien Supérieur en Administration de Bases de Données}{INSFP Hamdi Ben Yahia}{}{}{}
\cventry{2021}{Baccalauréat — Mathématiques}{Bouira}{}{}{}
\cventry{2023--2024}{Certifications IA (Computer Vision, NLP, LLM)}{Kaggle · Udemy · Hugging Face}{}{}{}

% LANGUES
\section{Langues}
\cvitem{Arabe}{C2 — Natif}
\cvitem{Français}{C1}
\cvitem{Anglais}{B2}

% PORTFOLIO
\section{Portfolio \& GitHub}
\cvitem{}{\url{https://ouahabi-benhenni.web.app} \quad \url{https://github.com/Greyma}}

\end{document}
